\chapter*{Preface}
\addcontentsline{toc}{chapter}{Preface}
\markboth{Preface}{Preface}

This book is a guide to \emph{FreeCAD}, the open-source parametric 3D
modeler, with a deliberate emphasis on its finite-element analysis (FEA)
capabilities and the mathematical theory that underpins them. It is
written for engineers and engineering students who already possess a
working knowledge of mechanics and calculus and who want both to
\emph{use} FreeCAD's FEM Workbench competently and to \emph{understand}
what the solver is actually computing.

\section*{How this book is organised}

The book is divided into three parts:

\begin{description}
  \item[Part~I --- FreeCAD Foundations.] A brisk, no-nonsense tour of
    the concepts a new user must internalise: the document object model,
    the workbench paradigm, the constraint-based Sketcher, and the two
    principal solid-modeling paradigms (\textsf{Part} and
    \textsf{PartDesign}). Readers already fluent in FreeCAD may skim
    this part and proceed to Part~II.
  \item[Part~II --- The Finite Element Method: Theory.] A self-contained,
    rigorous development of the FEM for linear elastostatics and its
    common extensions (modal, thermal, nonlinear, contact). We derive the
    strong and weak forms, the Galerkin discretisation, isoparametric
    elements, numerical quadrature, assembly, and the treatment of
    boundary conditions --- with equations, not hand-waving. This part
    stands on its own as a compact FEM primer.
  \item[Part~III --- The FEM Workbench in Practice.] The bridge between
    theory and software: how FreeCAD's FEM Workbench maps onto the
    machinery of Part~II, how the CalculiX and Elmer solvers are driven,
    a fully worked linear-static example, and how to post-process and
    \emph{verify} results.
\end{description}

\section*{On sources}

This guide is grounded in the official FreeCAD documentation, the
CalculiX and Elmer solver manuals, and the standard FEM literature
(Zienkiewicz \& Taylor, Hughes, Bathe, Cook \emph{et al.}). Software
behaviour is described for the FreeCAD~1.x generation; because FreeCAD is
under active development, always cross-check UI details against the live
documentation.

\vspace{1em}
\noindent\textit{Notation.} A consolidated notation reference is given in
Appendix~\ref{app:notation}. Vectors and matrices are set in bold; the
Voigt convention used throughout orders the stress components as
$[\sigma_{xx},\sigma_{yy},\sigma_{zz},\tau_{yz},\tau_{xz},\tau_{xy}]^{\mathsf T}$
and uses \emph{engineering} shear strains.
